\PaperPage{23}{3.1 The scalar resolvent and its universal domain}
\label{res:domain}
Throughout this chapter, $\sigma$ is a finite positive Borel measure on
$I=[0,1]$, $M=\sigma(I)$, and $\kappa>0$. Probability normalization is useful
for the application, but the analytic arguments require only finiteness.
The scalar object attached to these data is
\[
 S_\kappa(z)=\kappa\int_I\frac{1}{1+\kappa qz}\,d\sigma(q).
\]
The factor $\kappa$ and the plus sign in the denominator are part of our
convention. They determine both the Taylor coefficients and the direction
of the imaginary-part inequality. Changing either convention requires
changing the later formulas accordingly.

\begin{proposition}[Universal denominator domain]
The set on which every denominator is nonzero is exactly
\[
 D_\kappa:=\{z\in\mathbb C:\ 1+\kappa qz\ne0\text{ for all }q\in I\}
       =\mathbb C\setminus(-\infty,-1/\kappa].
\]
In particular, it contains the origin and both open half-planes.
\end{proposition}
\begin{proof}
At $q=0$ the denominator equals one. For $q>0$, its only zero is
\[
 z=-\frac{1}{\kappa q}\le-\frac1\kappa.
\]
Conversely, if $x\le-1/\kappa$, the number $q=-1/(\kappa x)$ belongs to
$(0,1]$ and makes that denominator zero. These two implications give the
claimed equality. For fixed $z\in D_\kappa$, the integrand is continuous
on the compact interval, hence bounded and integrable against $\sigma$.
\end{proof}

The adjective universal is essential. The displayed slit is determined by
all possible parameters in the interval, not by the support of one chosen
measure. For the point mass at $a\in(0,1]$, for example,
\[
 S_\kappa(z)=\frac{\kappa}{1+\kappa az},
 \qquad \sigma=\delta_a,
\]
which extends meromorphically through every other point of the slit.
For $\sigma=\delta_0$ the transform is the constant $\kappa$ and is entire.
Nevertheless, every excluded point is a genuine pole for some point mass.
Thus no larger common domain can guarantee a holomorphic transform for
all positive measures on $I$.

This distinction prevents a support bound from being mistaken for an
assertion about the exact singularities of the resulting function. We
will establish holomorphy on $D_\kappa$ without assuming the existence,
absence, or location of zeros of any entire function to which the
transform may later be compared.
All integrals here are ordinary complex integrals, equivalently the
integrals of the real and imaginary parts. No density with respect to
Lebesgue measure is required: discrete, continuous, and mixed measures
are treated by the same definition.

\PaperPage{24}{3.2 Compact denominator control and continuity}
\label{res:compact-control}
Pointwise nonvanishing alone does not justify an interchange of a limit
and an integral. What makes the present construction tractable is the
compactness of the parameter interval. It gives a positive lower bound
which persists throughout a neighborhood of the complex variable.

\begin{lemma}[A uniform local bound]
For every $z_0\in D_\kappa$, there are $c>0$ and $r>0$ such that
\[
 |1+\kappa qz|\ge c
 \quad(q\in I,\ |z-z_0|<r).
\]
\end{lemma}
\begin{proof}
The continuous positive function $q\mapsto|1+\kappa qz_0|$ attains a
strictly positive minimum $a$. Since $q\le1$, the reverse triangle
inequality gives
\[
 |1+\kappa qz|\ge a-\kappa|z-z_0|.
\]
Take $c=a/2$ and $r=a/(2\kappa)$. The same estimate also proves that this
ball lies in $D_\kappa$, so the domain is open.
\end{proof}

More generally, if $K$ is a nonempty compact subset of $D_\kappa$,
continuity on $K\times I$ gives
\[
 a_K:=\min_{z\in K,\,q\in I}|1+\kappa qz|>0,
 \qquad \sup_{z\in K}|S_\kappa(z)|\le\frac{\kappa M}{a_K}.
\]
The first assertion uses positivity at every pair: compactness then
excludes a sequence of denominators tending to zero. There is no need
to estimate the distance to each point of the slit separately.

For continuity, choose a neighborhood with the common bound $c$ and
subtract the two kernels before integrating. The elementary identity
\[
 \frac1{1+\kappa qz}-\frac1{1+\kappa qw}
 =\frac{\kappa q(w-z)}{(1+\kappa qz)(1+\kappa qw)}
\]
implies the local Lipschitz estimate
\[
 |S_\kappa(z)-S_\kappa(w)|\le\frac{\kappa^2M}{c^2}|z-w|.
\]
Here both points lie in the chosen neighborhood. This proves continuity
directly and supplies an explicit modulus there. Equivalently, dominated
convergence applies because the kernels are bounded by $c^{-1}$ and the
constant function is integrable for a finite measure.

The same mechanism will handle higher derivatives. Powers of the
coordinate remain bounded by one, while powers of the denominator are
controlled by this single lower bound. The resulting majorants depend
on the derivative order and the neighborhood, but not on the integration
variable. This uniformity is the substantive analytic input in the
compact-kernel argument formalized in the source development.
It also controls perturbations of the measure. If $\tau$ is another
finite positive measure and its transform is denoted $S_{\kappa,\tau}$,
integration against the total variation of the signed difference gives
\[
 \sup_{z\in K}|S_{\kappa,\sigma}(z)-S_{\kappa,\tau}(z)|
 \le\frac{\kappa}{a_K}\,|\sigma-\tau|(I).
\]
Thus convergence in total variation implies uniform convergence on
every compact subset of the common denominator domain.

\PaperPage{25}{3.3 Differentiation under the integral}
\label{res:differentiation}
Introduce the kernel tower and its integral by
\[
 K_n(z,q)=\frac{q^n}{(1+\kappa qz)^{n+1}},
 \qquad A_n(z)=\int_I K_n(z,q)\,d\sigma(q),\qquad n\ge0.
\]
For each fixed parameter, ordinary complex differentiation gives
\[
 \partial_zK_n(z,q)=-(n+1)\kappa K_{n+1}(z,q).
\]
The issue is to integrate this identity without relying on a purely
formal differentiation rule. Fix $z_0\in D_\kappa$ and a ball on which
the denominator has the lower bound $c>0$ from the preceding lemma.
On that ball the relevant estimates are
\[
 |K_n(z,q)|\le c^{-n-1},\qquad
 |\partial_zK_n(z,q)|\le(n+1)\kappa c^{-n-2}.
\]
Both majorants are integrable because $M<\infty$. Measurability follows
from continuity in $q$; it is not an additional assumption on the
measure beyond its Borel structure.

\begin{proposition}[Derivative of the kernel integral]
Every $A_n$ is holomorphic on $D_\kappa$, and
\[
 A_n'(z)=-(n+1)\kappa A_{n+1}(z).
\]
\end{proposition}
\begin{proof}
For nonzero complex $h$ sufficiently small, the segment from $z_0$ to
$z_0+h$ stays in the chosen ball. Apply the real fundamental theorem of
calculus to the smooth function of $t\in[0,1]$ obtained by restricting
the rational kernel to this segment. It yields
\[
 \frac{K_n(z_0+h,q)-K_n(z_0,q)}h
 =\int_0^1\partial_zK_n(z_0+th,q)\,dt.
\]
Consequently the difference quotient is bounded in modulus by the
constant derivative majorant, uniformly in $q$ and in small $h$.
Pointwise in $q$ it converges to $\partial_zK_n(z_0,q)$. Dominated
convergence therefore moves this limit through the $\sigma$ integral,
proving complex differentiability and the stated derivative formula at
$z_0$. Since $z_0$ was arbitrary and the domain is open, $A_n$ is
holomorphic there.
\end{proof}

This proof accommodates arbitrary directions of complex increments.
It does not invoke a real mean-value theorem for complex-valued
functions. It also explains why the derivative bound must hold on a
neighborhood, rather than merely at the point where differentiation is
performed. Repeating the proposition gives derivatives of every order
and, by the usual complex-analytic theorem, convergent local Taylor
expansions. These are the analytic germs needed for comparison with
an actual logarithmic derivative.
Atomic measures require no special exception: integration simply
includes their point evaluations, and the same constant bounds apply.
Nor does the proof need the measure to be a probability measure.
Its total mass enters only through the integrals of the majorants,
which are finite multiples of $M$.

\PaperPage{26}{3.4 The derivative tower and identification of germs}
\label{res:germ}
Write $u_n=\int_Iq^n\,d\sigma(q)$ for the moments of the given measure.
Since $S_\kappa=\kappa A_0$, induction using the preceding derivative
formula proves
\[
 S_\kappa^{(n)}(z)
 =(-1)^n n!\,\kappa^{n+1}
   \int_I\frac{q^n}{(1+\kappa qz)^{n+1}}\,d\sigma(q).
\]
Indeed, differentiation of the formula at order $n$ multiplies the
integral by $-(n+1)\kappa$ and increases its index by one, producing
exactly the coefficient at order $n+1$. Evaluation at the origin gives
\[
 S_\kappa^{(n)}(0)=(-1)^n n!\,\kappa^{n+1}u_n.
\]
There is also a direct convergence check: for $|\kappa z|<1$, the
geometric series of the kernel converges uniformly in $q\in I$.
Its terms are bounded by the summable series $|\kappa z|^n$, so
\[
 S_\kappa(z)=\sum_{n=0}^{\infty}(-1)^n\kappa^{n+1}u_nz^n,
 \qquad |z|<1/\kappa.
\]
This radius is a guaranteed lower bound; a particular measure may give
a larger radius.

\begin{proposition}[An analytic identification criterion]
If $R$ is holomorphic near zero and its derivatives there are
$R^{(n)}(0)=(-1)^n n!\kappa^{n+1}u_n$ for every $n$, then $R=S_\kappa$
on some neighborhood of zero.
\end{proposition}
\begin{proof}
Choose a disk on which the Taylor expansion of $R$ converges to $R$,
and intersect it with $|z|<1/\kappa$. Equality of the derivatives gives
equality term by term of the two convergent Taylor series on this
smaller disk, hence equality of their sums.
\end{proof}

For the actual application, set $\kappa=\texttt{actualKappa}$.
The fixed function $\Theta$ has
$\Theta(0)\ne0$, so $R=\Theta'/\Theta$ is holomorphic near zero. The
source's exact normalization and its retained moment-reality theorem
provide the identity
\[
 R^{(n)}(0)=(-1)^n n!\,\kappa^{n+1}u_n^{\mathrm{act}}.
\]
Thus a positive measure whose moments equal the actual normalized
moments identifies the actual germ. The representability assertion is
a mathematical hypothesis to be supplied by the preceding positivity
argument; it is not obtained by defining a new sequence from an
arbitrary measure. Likewise, agreement of finitely many moments gives
only finitely many matching derivatives. The conclusion here requires
all orders and analyticity, and therefore does not follow from a
truncated computation or a merely formal coefficient list.
For comparison, a finite truncation does have a rigorous error bound.
Since $0\le u_n\le M$, the geometric tail estimate yields
\[
 \left|S_\kappa(z)-\sum_{n=0}^N(-1)^n\kappa^{n+1}u_nz^n\right|
 \le\frac{\kappa M(\kappa|z|)^{N+1}}{1-\kappa|z|}.
\]
This bound controls the represented transform locally, but does not
replace the exact all-orders comparison with the prescribed function.

\PaperPage{27}{3.5 The imaginary-part identity in two representations}
\label{res:imaginary-part}
Let $z=x+iy$ with $y>0$. Conjugating the denominator gives the
pointwise identity
\[
 \operatorname{Im}\frac{\kappa}{1+\kappa qz}
 =-\frac{\kappa^2qy}{|1+\kappa qz|^2}.
\]
Taking imaginary parts commutes with integration because it is a
continuous real-linear map and the kernel is integrable. Consequently
\[
 \operatorname{Im}S_\kappa(z)
 =-\kappa^2y\int_I\frac{q}{|1+\kappa qz|^2}\,d\sigma(q)\le0.
\]
The sign follows from positivity of the measure and of the coordinate.
It is strict precisely when $\sigma((0,1])>0$: in that case the
nonnegative integrand is positive on a set of positive measure. A
measure concentrated at zero instead gives a real constant transform.
Thus the nonpositive formulation includes all admissible measures.

The Hilbert-space representation explains the same calculation
geometrically. On $H=L^2(I,\sigma;\mathbb C)$ take $T=M_q$ and $v=1$.
We use inner products conjugate-linear in the first argument and linear
in the second. Multiplication by the reciprocal denominator is bounded
on $D_\kappa$, giving
\[
 S_\kappa(z)=\kappa\langle v,(I+\kappa zT)^{-1}v\rangle,
 \qquad \|v\|^2=M.
\]
The following algebra applies equally to any bounded positive
selfadjoint contraction, whenever the displayed inverse exists.
Put $h=(I+\kappa zT)^{-1}v$; then $v=h+\kappa zTh$, and
\[
 \langle v,h\rangle
 =\|h\|^2+\kappa\overline z\langle Th,h\rangle,
 \qquad \langle Th,h\rangle=\langle h,Th\rangle\ge0.
\]
Multiplication by $\kappa$ and extraction of the imaginary part yield
\[
 \operatorname{Im}S_\kappa(z)
 =-\kappa^2y\langle h,Th\rangle.
\]
In the multiplication model, the quadratic form on the right is
exactly the integral in the scalar identity. The conjugate on $z$ is
forced by the inner-product convention; omitting it reverses the sign.

This is an anti-Herglotz convention: the upper half-plane is mapped
into the closed lower half-plane. A resolvent written with $I-zT$
instead has the opposite sign. The concrete operator calculation is
mathematical exposition of the measure model. Its validity does not
by itself assert that every corresponding Hilbert-space packaging is
an exported theorem of the accompanying Lean development.
There is a complementary symmetry identity:
\[
 S_\kappa(\overline z)=\overline{S_\kappa(z)}.
\]
Indeed, the parameters and the measure are real, so conjugation can
be moved through the integral. This gives the corresponding nonnegative
imaginary-part inequality in the lower half-plane. Both sign statements
are consequences of one positive representation, with no assumption
about zeros of a separate entire function.

\PaperPage{28}{3.6 Quantitative bounds and endpoint masses}
\label{res:bounds}
The denominator estimates give useful controls without prescribing the
singular behavior along the excluded slit. For $|z|\le r<1/\kappa$,
the reverse triangle inequality and the derivative tower give
\[
 |S_\kappa^{(n)}(z)|
 \le\frac{n!\,\kappa^{n+1}M}{(1-\kappa r)^{n+1}},\qquad n\ge0.
\]
On the nonnegative real axis all denominators are positive. In
particular, the transform is completely monotone there:
\[
 (-1)^nS_\kappa^{(n)}(x)
 =n!\,\kappa^{n+1}\int_I\frac{q^n}{(1+\kappa qx)^{n+1}}\,d\sigma(q)
 \ge0,\qquad x\ge0.
\]
For $n=0$ this also gives $S_\kappa(x)\le\kappa M$; if $M>0$ the
transform is strictly positive. The first derivative is nonpositive,
although strict decrease fails when all mass is at zero.

Let $a_0=\sigma(\{0\})$ and $a_1=\sigma(\{1\})$. Splitting the measure
into its endpoint atoms and its restriction to the open interval gives
the exact decomposition
\[
 S_\kappa(z)=\kappa a_0+\frac{\kappa a_1}{1+\kappa z}
       +\kappa\int_{(0,1)}\frac{d\sigma(q)}{1+\kappa qz}.
\]
At positive real infinity, the kernel is bounded by one and converges
pointwise to the indicator of $\{0\}$. Dominated convergence proves
\[
 \lim_{x\to+\infty}S_\kappa(x)=\kappa a_0.
\]
No convergence rate follows merely from finiteness: mass may accumulate
arbitrarily close to zero.

There is also a precise one-sided statement at the other endpoint.
For $0<\kappa\varepsilon\le1$, the quantity
$\kappa\varepsilon/(1-q+\kappa q\varepsilon)$ lies between zero and
one, because its denominator is at least $\kappa\varepsilon$.
It tends to zero for $q<1$ and equals one for $q=1$. Hence
\[
 \lim_{\varepsilon\downarrow0}
 \varepsilon S_\kappa(-1/\kappa+\varepsilon)=a_1.
\]
This follows by dominated convergence along real points in $D_\kappa$.
It identifies the endpoint atom without asserting meromorphic
continuation across that endpoint. If additional mass accumulates at
one, the remaining integral need not be holomorphic in any full
neighborhood of $-1/\kappa$. Accordingly, the exact atomic summand and
the proved one-sided limit must not be promoted to an unsupported
claim that the entire transform has an isolated pole there.
The disk bound is sharp uniformly over all measures: mass $M$ at
$q=1$ attains equality at the negative real point $z=-r$.
Likewise, the integral formula shows that every alternating derivative
of positive order is strictly positive on the nonnegative real axis
exactly when some mass lies away from zero. These observations identify
the degenerate cases without imposing additional regularity assumptions
on the measure.

\PaperPage{29}{3.7 Continuing an analytic defect from an origin germ}
\label{res:continuation}
An identity of logarithmic derivatives near zero must be propagated
without presupposing that the entire function has no zeros elsewhere.
The appropriate object for the identity theorem is a holomorphic
defect, defined without dividing by that function. Set
\[
 \mathbb H=\{z\in\mathbb C:\operatorname{Im}z>0\}.
\]
\begin{theorem}[Continuation of a differential identity]
Let $F$ be entire and let $S$ be holomorphic on $\mathbb H$.
Suppose that for some $\varepsilon>0$ one has
$F'(z)=S(z)F(z)$ whenever $z\in\mathbb H$ and $|z|<\varepsilon$.
Then this identity holds throughout $\mathbb H$.
\end{theorem}
\begin{proof}
The function
\[
 D(z)=F'(z)-S(z)F(z)
\]
is holomorphic on $\mathbb H$. Although zero is not in this open
half-plane, the hypothesized germ supplies an interior seed. Namely,
\[
 z_* = i\varepsilon/2\in\mathbb H\cap B(0,\varepsilon).
\]
A sufficiently small open disk about $z_*$ lies in this intersection,
and $D$ vanishes on that disk. The upper half-plane is convex, hence
connected. The identity theorem therefore gives $D=0$ on all of
$\mathbb H$, which is the required differential identity.
\end{proof}

To apply the theorem, suppose $F(0)\ne0$ and its logarithmic derivative
agrees near zero with a function $S$ holomorphic both there and on
$\mathbb H$. Continuity allows the germ neighborhood to be shrunk so
that $F$ has no zeros in it. Multiplication by $F$ then gives
\[
 \frac{F'}F=S\quad\hbox{near }0
 \quad\Longrightarrow\quad F'=SF\quad\hbox{on }\mathbb H.
\]
Only after continuation do we divide at individual points where the
denominator is known to be nonzero:
\[
 F(z)\ne0,\ z\in\mathbb H
 \quad\Longrightarrow\quad \frac{F'(z)}{F(z)}=S(z).
\]
Thus any imaginary-part bound for $S$ transfers to the logarithmic
derivative at precisely the points needed for pole exclusion.

For the measure resolvent, holomorphy on $\mathbb H$ follows from
$\mathbb H\subset D_\kappa$. For the actual entire function, the
previous germ criterion supplies the local equality when its moments
are represented by $\sigma$. The argument keeps these two inputs
separate: positivity constructs and controls $S$, while analyticity
and exact moment agreement identify the germ. Global zero-freeness is
not one of the inputs to this continuation step.
The defect identity itself remains meaningful at a hypothetical zero.
At such a point it states $F'=0$, which by itself would only exclude
a simple zero. The stronger conclusion comes from the identity on a
whole neighborhood and the local behavior of logarithmic derivatives,
examined next.

\PaperPage{30}{3.8 The pole probe and conjugate-zero symmetry}
\label{res:pole-probe}
The sign constraint has a local consequence for any nontrivial entire
function. Its proof uses the positive integral multiplicity of a
holomorphic zero, independently of how the sign constraint was obtained.

\begin{theorem}[Upper-half-plane pole obstruction]
Let $F$ be entire and not identically zero. If
\[
 \operatorname{Im}\frac{F'(z)}{F(z)}\le0
 \quad\text{for every }z\in\mathbb H\text{ with }F(z)\ne0,
\]
then $F$ has no zeros in $\mathbb H$.
\end{theorem}
\begin{proof}
Suppose $\rho\in\mathbb H$ is a zero. Local factorization gives a
positive integer $m$, a disk about $\rho$, and a holomorphic function
$g$ without zeros on that disk such that
\[
 F(z)=(z-\rho)^m g(z),\qquad
 \frac{F'(z)}{F(z)}=\frac{m}{z-\rho}+B(z),\quad B=\frac{g'}g,
\]
where the quotient identity holds away from $\rho$. On a smaller
closed disk, continuity of $B$ gives a finite bound
$|\operatorname{Im}B|\le K$ with $K\ge0$. Choose $\varepsilon>0$
small enough that $\rho-i\varepsilon$ lies in this smaller disk,
$\varepsilon<\operatorname{Im}\rho$, and $\varepsilon<m/(K+1)$.
The probe point is in $\mathbb H$ and is not a zero, because the
factor $g$ does not vanish. At that point,
\[
 \operatorname{Im}\frac{F'(\rho-i\varepsilon)}{F(\rho-i\varepsilon)}
 =\frac{m}{\varepsilon}+\operatorname{Im}B(\rho-i\varepsilon)
 \ge\frac{m}{\varepsilon}-K>0.
\]
This contradicts the assumed sign inequality.
\end{proof}

The downward direction is decisive: the imaginary part of
$m/(-i\varepsilon)$ is positive. An upward probe would give the
opposite sign and would not establish this contradiction. The local
factorization also ensures that the tested quotient is defined;
there is no appeal to a value at the putative pole itself.

If the zeros are invariant under conjugation, the lower half-plane
is excluded as well. Indeed,
\[
 F(\rho)=0\Longrightarrow F(\overline\rho)=0,
 \qquad \operatorname{Im}\rho<0
 \Longrightarrow\operatorname{Im}\overline\rho>0.
\]
A sufficient condition is Schwarz symmetry
$F(\overline z)=\overline{F(z)}$, although only the implication for
zeros is needed. Combining this symmetry with the theorem leaves
real zeros as the only possibility. Applied after the analytic-defect
continuation, this is the general mechanism by which a positive
measure representation constrains the zeros of the actual entire
function. The next chapter specifies the actual-function bindings
needed to turn that constraint into the repository's stated endpoint.
The limitation to nonreal zeros is necessary. For example, if $a$ is
real, $F(z)=z-a$ satisfies the sign hypothesis because
\[
 \operatorname{Im}\frac1{z-a}
 =-\frac{\operatorname{Im}z}{|z-a|^2}<0\qquad(z\in\mathbb H),
\]
yet has a real zero at $a$. The theorem therefore imposes exactly
the stated location constraint.
